\section{March 06}

\subsection{Frobenius element (continue)}

    \begin{proposition}
        We have a tower of fields \(M/L/K\).
        Suppose that \(M/K\) is Galois.
        Let \(D \in \calO_M\) a prime ideal and \(\frakP = D \cap \calO_L\) and \(\frakp = D \cap \calO_K\).
        Assume that \(L/K\) is Galois, then \(\Frob_{(D, M/K)}|_L = \Frob_{(\frakP, L/K)}\).
    \end{proposition}

    \begin{corollary}\label{cor:split_completely_and_frobenius}
        Let \(L/K\) be a Galois extension and \(\frakp\) a prime of \(K\) unramified in \(L\).
        Then 
        \(\frakp\) is split completely in \(L\) if and only if \(\Frob_{(\frakP, L/K)} = 1\) for one (or equivalently, any) prime \(\frakP\) of \(L\) above \(\frakp\).
    \end{corollary}

    \begin{proposition}\label{prop:frobenius_in_compositum}
        Let \(L_1, L_2\) be Galois extensions of \(K\) and \(M=L_1L_2\).
        Then \(M/K\) is Galois and there is a natural embedding \(j:\Gal(M/K) \hookrightarrow \Gal(L_1/K) \times \Gal(L_2/K)\).
        Under the embedding \(j\), we have \(j(\Frob_{(D, M/K)}) = (\Frob_{(\frakP, L_1/K)}, \Frob_{(\frakP, L_2/K)})\).
    \end{proposition}

    Hence in the situation of \cref{prop:frobenius_in_compositum}, \(\frakp\) is split completely in \(M\) if and only if \(\frakp\) is split completely in both \(L_1\) and \(L_2\).

    \begin{example}
        Let \(K = \bbQ[\xi_n]\) with \(\xi_n = e^{2\pi i/n}\) and \(p\) an odd prime.
        If \(p \mid n\), then \(p\) is ramified in \(K\).
        The Frobenius element \(\Frob_{(p, K/\bbQ)}\) is not defined.
        Suppose that \(p \nmid n\).
        Then \(\sigma = \Frob_{(p, K/\bbQ)}\) is the unique element in \(\Gal(K/\bbQ) \cong (\bbZ/n\bbZ)^\times\) 
        such that \(\sigma(\alpha) = \alpha^p \pmod{\frakp}\) for all \(\alpha \in \calO_K\) and all prime \(\frakp\) of \(K\) above \(p\).
        We claim that \(\sigma\) is given by \(\sigma(\xi_n) = \xi_n^p\).
        Fix a prime \(\frakp\) of \(K\) above \(p\).
        We have 
        \[ \sigma(\sum a_i \xi_n^i) = \sum a_i \sigma(\xi_n)^i = \sum a_i \xi_n^{pi} \equiv \sum a_i^p \xi_n^{pi} \equiv \left(\sum a_i \xi_n^i\right)^p \pmod{\frakp}. \]
        Note that \(\sigma = \Frob_{(p, K/\bbQ)}\) has order \(f\) where \(f\) is the smallest positive integer such that \(p^f \equiv 1 \pmod{n}\).
    \end{example}

    \begin{example}
        Let \(K = \bbQ[\sqrt{d}]\) where \(d\) is a square-free integer.
        Let \(p\) be a prime unramified in \(K\).
        Identify \(\Gal(K/\bbQ)\) with \(\{\pm 1\}\) where \(-1\) acts on \(K\) by sending \(\sqrt{d}\) to \(-\sqrt{d}\).
        Then by \cref{cor:split_completely_and_frobenius} we have
        \[ \Frob_{(p, K/\bbQ)} = \begin{cases}
            1 & \text{if } d \text{ is a quadratic residue modulo } p, \\
            -1 & \text{if } d \text{ is not a quadratic residue modulo } p.
        \end{cases} \]
    \end{example}

    \paragraph{The quadratic reciprocity law by Frobenius elements}
    Let \(K = \bbQ[\xi]\) with \(\xi = \exp({2\pi i/p})\) where \(p\) is an odd prime.
    Then \(\Gal(K/\bbQ) \cong (\bbZ/p\bbZ)^\times\) is cyclic of order \(p-1\).
    There is a unique subgroup of \(\Gal(K/\bbQ)\) of index \(2\), 
    consisting of the squares in \((\bbZ/p\bbZ)^\times\),
    which corresponds to a unique quadratic extension \(F\) of \(\bbQ\) contained in \(K\).

    We check the ramification in \(K\) and \(F\).
    The only prime ramified in \(K\) is \(p\).
    Hence the only prime ramified in \(F\) is \(p\).
    If \(p \equiv 1 \pmod{4}\), then \(p\) is the only prime ramified in \(\bbQ(\sqrt{p})\).
    And \(\bbQ(\sqrt{p})\) is the unique quadratic extension of \(\bbQ\) ramified only at \(p\).
    Hence \(F = \bbQ(\sqrt{p})\).
    Similarly, if \(p \equiv 3 \pmod{4}\), then \(F = \bbQ(\sqrt{-p})\).
    Denote by \(d = (-1)^{\frac{p-1}{2}}p\), then \(F = \bbQ(\sqrt{d})\).

    Let \(q\) be an odd prime different from \(p\).
    Then \(q\) is unramified in \(F\).
    Consider the Frobenius element \(\Frob_{(q, K/\bbQ)}: \xi \mapsto \xi^q\).
    We have 
    \begin{align*}
            & q \text{ is a quadratic residue modulo } p \\
        \iff & \Frob_{(q, K/\bbQ)} \in \Gal(K/\bbQ) \text{is a square} \\
        \iff & \Frob_{(q, F/\bbQ)} = \Frob_{(q, K/\bbQ)}|_F \in \Gal(F/\bbQ) \text{is the identity} \\
        \iff & \text{the prime } q \text{ is split completely in } F \\
        \iff & d \text{ is a quadratic residue modulo } q
    \end{align*}
    Hence
    \[ \left(\frac{d}{q}\right) = \left(\frac{(-1)^{\frac{p-1}{2}}p}{q}\right) = \left(\frac{-1}{q}\right)^{\frac{p-1}{2}} \left(\frac{p}{q}\right) = \left(\frac{q}{p}\right). \]
    This gives a proof of the part of the quadratic reciprocity law for odd primes.

    \begin{theorem}[Quadratic reciprocity law]\label{thm:quadratic_reciprocity_law}
        Let \(p\) and \(q\) be two distinct odd primes.
        Then we have the following relations.
        \begin{enumerate}
            \item \(\left(\frac{p}{q}\right) \left(\frac{q}{p}\right) = (-1)^{\frac{p-1}{2} \cdot \frac{q-1}{2}}\).
                In other words, 
                \[\left(\frac{p}{q}\right) = \begin{cases}
                    \left(\frac{q}{p}\right) & \text{if } p \equiv 1 \text{ or } q \equiv 1 \pmod{4}, \\
                    -\left(\frac{q}{p}\right) & \text{if } p \equiv q \equiv 3 \pmod{4}.
                \end{cases} \]
            \item \(\left(\frac{-1}{p}\right) = (-1)^{\frac{p-1}{2}}\).
            \item \(\left(\frac{2}{p}\right) = (-1)^{\frac{p^2-1}{8}}\).
        \end{enumerate}
    \end{theorem}
    We finish the proof of quadratic reciprocity law here.
    \begin{proof}
        The part (b) follows from the Fermat's theorem on sums of two squares.
        Then the part (a) follows from the discussion above.

        For the part (c), let \(\xi\) be a primitive \(8\)-th root of unity in \(\overline{\bbF_p}\).
        Set \(a = \xi + \xi^{-1}\).
        From \(\xi^4 = -1\) we deduce that
        \[ X^4+1 = (X^2 - \xi^2)(X^2 - \xi^{-2}) \in \overline{\bbF_p}[X]. \]
        This implies that \(\xi^2 + \xi^{-2} = 0\) and hence \(a^2 = 2\).
        Note that the equality
        \[ \xi^p + \xi^{-p} = \xi + \xi^{-1} \]
        holds if and only if \(p \equiv \pm 1 \pmod{8}\).
        Since \(\xi^p + \xi^{-p} = a^p\), we have 
        \[ a^p = a \iff p \equiv \pm 1 \pmod{8} \iff a \in \bbF_p. \]
        We are done.
    \end{proof}


\subsection{The Chebotarev density theorem}

    \begin{definition}
        Let \(K\) be a number field and \(P= \mSpec \calO_K\) the set of finite places of \(K\).
        For every \(S \subset P\), we say that \(S\) has \emph{natrual density} \(\delta\) if
        \[ \lim_{N \to \infty} \frac{\#\{\frakp \in S: \Nm(\frakp) \leq N\}}{\#\{\frakp \in P: \Nm(\frakp) \leq N\}} = \delta. \]
    \end{definition}

    \begin{theorem}[Chebotarev density theorem]\label{thm:Chebotarev_density_theorem}
        Let \(L/K\) be a Galois extension of number fields with Galois group \(G = \Gal(L/K)\).
        Let \(C\) be a conjugacy class in \(G\).
        Then the set of primes \(\frakp\) of \(K\) such that \(\Frob_{(\frakp, L/K)} = C\) has natural density \(\frac{|C|}{|G|}\).
    \end{theorem}

    Proof of \cref{thm:Chebotarev_density_theorem} will be given later.

    \begin{corollary}
        Let \(L/K\) be a Galois extension of number fields.
        The set of primes of \(K\) that split completely in \(L\) has natural density \(\frac{1}{[L:K]}\).
    \end{corollary}

    \begin{remark}
        There exists a bound for the error in terms of discriminant of the polynomial, 
        but it is large.
        Given a polynomial \(f \in \bbZ[X]\), there exists a bound \(B\) 
        such that if a given cycles type does not appear as the Frobenius element for any prime \(p \leq B\), 
        then it does not appear as the Frobenius element for any prime \(p\).

        Here, by ``cycle type'' we mean the cycle type of the Frobenius element as a permutation of the roots of \(f \pmod{p}\).
    \end{remark}

    \begin{example}
        Let \(K = \bbQ[\xi_n]\), \(\Gal(K/\bbQ) \cong (\bbZ/n\bbZ)^\times\).
        We know that \(\Frob_{(p, K/\bbQ)}\) corresponds to \(p \in (\bbZ/n\bbZ)^\times\) for any prime \(p\) unramified in \(K\).
        The Chebotarev density theorem implies that 
        the set of primes \(p\) such that \(p \equiv a \pmod{n}\) are equi-distributed of density \(\frac{1}{\varphi(n)}\) for each \(a \in (\bbZ/n\bbZ)^\times\).
        This implies the Dirichlet's theorem on arithmetic progressions.
    \end{example}

    \paragraph{Application}

    Let \(L/K\) be a Galois extension of number fields with Galois group \(G\).
    Denote by 
    \[ \upoperatorname{Spl}(L/K) = \{\frakp \in \mSpec O_K: \frakp \text{ is split completely in } L\} \]
    the set of primes of \(K\) that split completely in \(L\).
    
    \begin{theorem}\label{thm:extension_and_splitting_primes}
        If \(L/K\) and \(M/K\) is Galois extensions of number fields.
        Then \(L \subset M\) if and only if \(\upoperatorname{Spl}(M/K) \subset \upoperatorname{Spl}(L/K)\).
    \end{theorem}
    \begin{proof}
        (\(\Rightarrow\)) obvious.

        (\(\Leftarrow\)) We have 
        \[ \upoperatorname{Spl}(LM/K) = \upoperatorname{Spl}(L/K) \cap \upoperatorname{Spl}(M/K) = \upoperatorname{Spl}(M/K) \]
        by \cref{prop:frobenius_in_compositum}.
        By the Chebotarev density theorem, we know that 
        \[ \frac{1}{[LM:K]} = \text{ density of } \upoperatorname{Spl}(LM/K) = \text{ density of } \upoperatorname{Spl}(M/K) = \frac{1}{[M:K]}. \]
        % density of \(\upoperatorname{Spl}(LM/K)\) is \([LM:M]\) times the density of \(\upoperatorname{Spl}(M/K)\).\Yang{}
        Hence \(LM = M\) and then \(L \subset M\).
    \end{proof}

    \begin{corollary}
        Same as \cref{thm:extension_and_splitting_primes}.
        Then 
        \[ L = M \iff \upoperatorname{Spl}(L/K) = \upoperatorname{Spl}(M/K). \]
    \end{corollary}

    \begin{slogan}
        We can use the information on base field (e.g. the splitting of primes) to determine the extension field.
    \end{slogan}

    \begin{theorem}
        For any number field \(K\), integer \(N\) and finite set \(S\) of primes of \(K\), 
        there are only finite many fields extensions \(L/K\) of degree \(N\) unramified outside \(S\).
    \end{theorem}
    \begin{remark}
        Given a compact Riemann surface \(X\) and \(P_1, \dots, P_n\) points on \(X\), 
        there are only finite many \'etale cover of \(X\setminus\{P_1, \dots, P_n\}\) of degree \(N\).
        This is since \(\pi_1(X\setminus\{P_1, \dots, P_n\})\) is finitely generated.

        For a finitely generated group \(G\) and integer \(N\), there are only finite many subgroups of \(G\) of index \(N\).
        This is since the subgroups of \(G\) of index \(N\) corresponds a homomorphism \(G \to S_N\) by \(G \acts G/H\) and there are only finite many homomorphisms \(G \to S_N\).
    \end{remark}
    \begin{proof}
        Recall for all prime \(v\) and integer \(N\), there are only finite many extensions of \(K_v\) of degree dividing \(N\).
        This since every extension of \(K_v\) can be decomposed into an unramified extension and a totally ramified extension.
        For the unramified extension, there are only finite many of them since they are determined by the residue field extension.
        For the totally ramified extension, they are determined by the Eisenstein polynomials.
        And there exists \(M\) depending only on \(N\) such that for two Eisenstein polynomials $f, g$ of degree \(N\) with coefficients in \(\calO_{K_v}\),
        \[ f \equiv g \pmod{\frakm_v^M} \implies f \text{ and } g \text{ define the same extension of } K_v. \]
        Hence the assertion follows from the finiteness of \(\calO_{K_v}/\frakm_v^M\).

        Next \(\disc(L/K) = \prod_{w|v} \disc(L_w/K_v) \in \calO_{K_v}\).
        Let \(L\) be an extension of \(K\) of degree \(N\) unramified outside \(S\).
        Then \(\disc(L/K)\) is bounded.
        The result follows from the Hermite's theorem below (\cref{thm:hermite_fix_discriminant})
    \end{proof}


    \begin{theorem}[Hermite]\label{thm:hermite_fix_discriminant}
        Let \(K\) be a number field. 
        There are only finite many extensions of \(K\) with given discriminant.
    \end{theorem}
    \begin{proof}
        For any extension \(L/K\), we have \(\disc(L/K) = \Nm_{K/\bbQ}(\disc(L/K)) \cdot \disc_K^{[L:K]}\).
        Hence we can assume that \(K = \bbQ\) and we want to show that there are only finite many extensions \(L\) of \(\bbQ\) with given discriminant \(d\).
        By Minkowski's theorem, we have 
        \[ |d| > \left(\frac{\pi}{4}\right)^{n/2} \frac{n^n}{n!}. \]
        Hence we may fix \(n = [L:\bbQ]\).
        Set \(D = |d|\).

        Let \(\sigma_1, \dots, \sigma_r \in M_{L,\infty}\) be the infinite places of \(L\) with completion \(\bbR\) and 
        \(\sigma_{r+1}, \dots, \sigma_{r+s}\) be the infinite places of \(L\) with completion \(\bbC\).
        Consider \(\Sigma: L \to \bbR^r \times \bbC^s\) given by \(x \mapsto (\sigma_1(x), \dots, \sigma_{r+s}(x))\).
        Define 
        \[ X \coloneqq \begin{cases}
            \{ (x_1, \dots, x_{r+s}) : |x_1| < \sqrt{D+1}, \quad |x_i| < 1, i = 2, \dots, r+s \} & \text{if } r \neq 0, \\
            \{ (x_1, \dots, x_s) : |\Real x_1|<1,|\Image x_1|<\sqrt{D+1},\quad |x_i| < 1, i = 2, \dots, s \} & \text{if } r = 0.
        \end{cases} \]
        We can compute the volume of \(X\):
        \[ \vol(X) = \begin{cases}
            2^r \pi^s \sqrt{D+1} & \text{if } r \neq 0, \\
            2 \pi^{s-1} \sqrt{D+1} & \text{if } r = 0.
        \end{cases} \]
        In either case, we have \(\vol(X) > 2^r \sqrt{D}\).
        By the Minkowski's theorem on the lattice \(\calO_L\) and the convex body \(X\), 
        there exists a nonzero element \(\alpha \in \calO_L\) such that \(\Sigma(\alpha) \in X\).
        All conjugates of \(\alpha\) are absolutely bounded by a constant depending only on \(D\).
        Hence there are only finite many possibilities of such \(\alpha \in \calO_{\overline{\bbQ}}\).

        By the construction of \(X\), if \(\sigma(\alpha) = \alpha\) for some \(\sigma \in \Gal(L/\bbQ)\), 
        then \(\sigma\) must be the identity on \(L\). 
        Hence \(L = \bbQ(\alpha)\).
    \end{proof}

